% Le MVCC de PostgreSQL : chaque UPDATE crée une version ; chaque transaction
% voit la version valide pour son instantané ; VACUUM retire ce que plus
% personne ne peut voir (valeurs de ctid et xmin mesurées par ch18/mvcc.py).
\documentclass[border=6pt]{standalone}
\usepackage{sd}
\begin{document}
\begin{tikzpicture}[x=1mm, y=1mm, note/.append style={font=\sffamily\small},
    v/.style={draw=#1, fill=#1Bg, line width=0.8pt, rounded corners=2pt, minimum width=34mm, minimum height=13mm, font=\sffamily\small, align=center, inner sep=2pt}]
  \node[v=Dslate] (v0) at (0,0) {v = 0\\{\footnotesize ctid (0,1), xmin 756, xmax 757}};
  \node[v=Dslate] (v1) at (58,0) {v = 1\\{\footnotesize ctid (442,109), xmin 757, xmax 758}};
  \node[v=Dgreen] (v2) at (116,0) {v = 2\\{\footnotesize ctid (442,110), xmin 758}};
  \draw[arr=Dslate] (v0) -- node[lbl, above] {UPDATE} (v1);
  \draw[arr=Dslate] (v1) -- node[lbl, above] {UPDATE} (v2);
  \node[pillo=Dblue] (t1) at (25,-22) {transaction ouverte avant 757 : voit v = 0};
  \node[pillo=Dgreen] (t2) at (110,-22) {nouvelle transaction : voit v = 2};
  \draw[arr=Dblue, dash pattern=on 2pt off 1.5pt] (t1.north) -- (v0.south);
  \draw[arr=Dgreen, dash pattern=on 2pt off 1.5pt] (t2.north) -- (v2.south);
  \node[note, anchor=north, align=center] at (58,-30) {v = 0 et v = 1 sont des versions mortes : VACUUM les retire quand plus aucune\\transaction ouverte ne peut les voir ; tant que la transaction bleue reste ouverte, il ne peut pas};
\end{tikzpicture}
\end{document}
